\documentclass{situationist-preprint}

\title{The Rationality of Self-Sacrifice}
\msid{ETH-25-a}
\msfamily{ETH-25}
\msdiscipline{Normative and Applied Ethics}
\mstype{Principal manuscript}
\mskeywords{self-sacrifice; rationality; well-being; egoism; caring; deepest concerns; Overvold's problem}

\begin{document}
\maketitle

\begin{abstract}
A self-sacrifice is an act that the agent expects to leave her worse off than an available alternative, done for the sake of something else. Egoists hold that such acts are always irrational; desire-satisfaction theorists of well-being have been charged with making them impossible, since an act done for what one most wants satisfies one's desires. We argue that self-sacrifice is possible, and that it is rational if and only if the agent's deepest concerns lie beyond her own well-being. The argument distinguishes well-being, the good of the agent's life for her, from the objects of her concerns, which may include her children, her work and her principles. Rationality, in the sense at issue, requires coherence between an agent's actions and her deepest concerns. If some of those concerns lie beyond her well-being, some sacrifices serve them and are rational; if none do, every sacrifice runs against all of them and is irrational. We show that the result is consistent with the existence of supererogatory sacrifice and with the thesis that rational agents are moved by their moral judgments, and we distinguish the rationality of a sacrifice from the reasons there may be for it.
\end{abstract}

\section{The problem}

Soldiers fall on grenades, parents forgo comforts for their children, and dissidents accept prison rather than recant. Each acts in a way she expects to leave her worse off, for the sake of something other than herself. The question whether such acts are rational has two sources.

The first is egoism. On rational egoism, an act is rational for an agent if and only if it maximizes her own good, so that self-sacrifice is irrational by definition \citep{Sidgwick1874,Parfit1984}. Sidgwick took the conflict between egoistic and impartial reasons to be irresolvable, and called it the dualism of practical reason; Parfit argued against the self-interest theory that it is not the only rational aim to make one's own life go best.

The second source is a puzzle about well-being, set out by \citet{Overvold1980}. If well-being consists in the satisfaction of one's desires, then an agent who does what she most wants, all things considered, makes her life go as well as possible. The soldier who falls on the grenade wants to save his comrades more than he wants to live; on the unrestricted desire theory he therefore maximizes his well-being, and his act is no sacrifice. Since self-sacrifice plainly occurs, the unrestricted desire theory is false, or the concept of self-sacrifice is confused. \citet{Heathwood2011} and \citet{Portmore2007} have discussed the puzzle, and their proposals restrict either the desires that count towards well-being or the notion of sacrifice.

We defend a biconditional. Self-sacrifice is rational if and only if the agent's deepest concerns lie beyond her own well-being. The left-to-right direction says that an agent whose concerns all terminate in her own good cannot rationally sacrifice it; the right-to-left direction says that an agent who cares about something other than her own good can.

\section{Framework}

\begin{definition}[Well-being]\label{d:wb}
An agent's \emph{well-being} is how well her life goes for her. It includes her experiences, her contact with reality, and the satisfaction of those of her desires whose satisfaction entails that she exists at the time of satisfaction.
\end{definition}

Definition~\ref{d:wb} adopts Overvold's restriction on the desire component, together with the result of ETH-05 that well-being includes contact with reality and is not confined to pleasure. It is neutral between the main theories of well-being in all other respects. Its effect is that the satisfaction of a desire that one's comrades survive, after one's death, does not contribute to one's well-being.

\begin{definition}[Self-sacrifice]\label{d:sac}
An act $a$ is a \emph{self-sacrifice} if the agent performs it for the sake of some object $o$, and she expects that $a$ will leave her well-being lower than some available alternative would.
\end{definition}

\begin{definition}[Deepest concerns]\label{d:deep}
An agent's \emph{deepest concerns} are the objects of her final caring that she would not give up on reflection and that structure her other concerns: the things whose flourishing she treats as reasons that her other desires answer to \citep{Frankfurt1982}. A concern \emph{lies beyond} her well-being if its object is not her well-being or a component of it.
\end{definition}

\begin{principle}[Rationality as coherence]\label{p:rational}
An act is rational for an agent if and only if, given her beliefs, no available alternative is better supported by the ordering of options that her deepest concerns jointly determine.
\end{principle}

Principle~\ref{p:rational} expresses rationality in the structural sense: the coherence of an agent's choices with what she most cares about. It is the sense in which the question about self-sacrifice is usually asked, namely whether the sacrificing agent makes a mistake by her own lights. It is distinct from the question whether she has reason to act as she does. ACT-12 argues that there are external reasons, which an agent may have although no sound deliberation from her motivations would reveal them. If so, an agent may have reason to sacrifice herself even when doing so would not be rational for her in the present sense, and conversely; our thesis concerns rationality only.

\section{The argument}

\begin{lemma}\label{l:possible}
Self-sacrifice is possible.
\end{lemma}

\begin{proof}
By Definition~\ref{d:wb}, the satisfaction of a desire whose object lies outside the agent's life does not contribute to her well-being. An agent can therefore expect an act to secure an object she cares about, such as her comrades' survival, while expecting it to lower her well-being. Such an act is a self-sacrifice by Definition~\ref{d:sac}. Overvold's puzzle arises only on the unrestricted desire theory, which Definition~\ref{d:wb} does not adopt.
\end{proof}

\begin{lemma}\label{l:if}
If some of an agent's deepest concerns lie beyond her well-being, then some self-sacrifice is rational for her.
\end{lemma}

\begin{proof}
Let $o$ be a deepest concern that lies beyond her well-being. There are possible situations in which serving $o$ requires a cost to her well-being: situations in which the only way to protect a child, a cause or a principle she cares about is to forgo some good of her own. In such a situation, if the importance she attaches to $o$, relative to her other deepest concerns, outweighs the cost, the ordering her concerns determine ranks the sacrifice above the alternatives, and by Principle~\ref{p:rational} the sacrifice is rational.
\end{proof}

\begin{lemma}\label{l:onlyif}
If none of an agent's deepest concerns lies beyond her well-being, then no self-sacrifice is rational for her.
\end{lemma}

\begin{proof}
Suppose all her deepest concerns are her well-being or its components. Then the ordering they determine ranks options by their expected contribution to her well-being. By Definition~\ref{d:sac}, a self-sacrifice is an act she expects to leave her well-being lower than some alternative; that alternative is ranked above it, and by Principle~\ref{p:rational} the sacrifice is not rational. Her caring about some further object $o$, if it is not among her deepest concerns, does not affect the ordering, since by Definition~\ref{d:deep} it answers to her deepest concerns and gives no independent support.
\end{proof}

\begin{theorem}\label{t:main}
Self-sacrifice is rational for an agent if and only if her deepest concerns lie, at least in part, beyond her own well-being.
\end{theorem}

\begin{proof}
By Lemma~\ref{l:possible}, the question is not empty. The right-to-left direction is Lemma~\ref{l:if}, and the left-to-right direction is the contrapositive of Lemma~\ref{l:onlyif}.
\end{proof}

The theorem is to be read with its quantifiers in view. On the right-to-left side, some self-sacrifice is rational: not every one. A sacrifice that does little or nothing for the agent's concerns, or more than they require, is dispreferred by the ordering they determine, and is irrational even for an agent whose concerns lie beyond herself. Self-abasement and pointless martyrdom are of this kind.

\section{Consistency with neighboring results}

\paragraph{Supererogation.} ETH-14 argues that supererogatory acts exist, bounded below by decency and above by sainthood, beyond which further sacrifice ceases to be morally better and becomes waste or self-abasement. Theorem~\ref{t:main} explains why the agent who performs a supererogatory sacrifice need not be irrational: her deepest concerns include the good of those she helps. It also explains why sacrifice beyond the saintly bound tends to be irrational as well as not admirable: a sacrifice disproportionate to the good it does is dispreferred by any ordering in which the agent's own well-being has some weight, as it has for almost everyone. The moral bound and the rational bound coincide when the agent's concerns are the moral ones.

\paragraph{Moral motivation.} MTE-07 argues that moral judgment necessarily motivates rational agents, and that amoralists, though possible, are practically irrational. Combined with Theorem~\ref{t:main}, this has a consequence. A rational agent who judges that she is required to make a sacrifice is moved by the judgment, so that a concern for what morality requires is among the concerns her choices answer to, and that concern lies beyond her well-being. Required sacrifices, such as those the duties of rescue and beneficence demand, are therefore rational for the rational agent who accepts the requirement.

\paragraph{The good.} ETH-01 lists well-being as one of seven irreducible components of the good, alongside knowledge, achievement, relationships, fairness, beauty and autonomy. An agent whose deepest concerns are directed to any of the other six, for their own sake and not as parts of her life going well for her, has concerns that lie beyond her well-being.

\section{Objections and replies}

\begin{objection}[Hidden egoism]\label{o:egoism}
Every agent who sacrifices herself does what she most wants, and so serves herself after all. The soldier would be unable to live with himself if he did not act.
\end{objection}

\begin{reply}
That the soldier does what he most wants does not show that he serves his well-being, unless well-being is the satisfaction of whatever one most wants, which Lemma~\ref{l:possible} rejects. Some sacrificers would indeed be unable to live with themselves; but the soldier who falls on a grenade does not live at all, and his act cannot be explained by the avoidance of later guilt. The objection describes an egoistic motive that some sacrificers may have; it does not show that all do.
\end{reply}

\begin{objection}[Self-interest theory]\label{o:si}
Rationality requires each agent to make her own life go as well as possible. Theorem~\ref{t:main} assumes a structural notion of rationality that begs the question against this theory.
\end{objection}

\begin{reply}
The self-interest theory is the claim that every agent has most reason to maximize her own good. \citet{Parfit1984} argued that the theory is unstable: it gives the agent's future self a special status that it cannot justify while denying any such status to other people. Whatever one makes of that argument, the self-interest theory is a claim about reasons, and we have separated the question of reasons from the question of rationality in the structural sense. If the self-interest theory is true, then no one has most reason to sacrifice herself; it would still be true that a sacrifice coheres with the concerns of an agent who cares about others, and that the agent does not err by her own lights.
\end{reply}

\begin{objection}[Ill-chosen concerns]\label{o:bad}
An agent may have deepest concerns that are worthless or vicious, such as devotion to a cult. Theorem~\ref{t:main} implies that her sacrifice for it is rational, which seems wrong.
\end{objection}

\begin{reply}
The theorem implies that her sacrifice coheres with what she cares about, and we accept this. The defect lies in the concerns, not in their relation to her action, and it is a defect that rationality in the structural sense does not measure. The cult member's error is a failure to respond to the reasons that the worth of things gives, and if there are external reasons, as ACT-12 argues, she has reasons she does not respond to. To call her sacrifice irrational would be to run together two ways of going wrong that should be kept apart.
\end{reply}

\section{Conclusion}

Self-sacrifice is possible, because well-being does not include the satisfaction of every concern, and it is rational if and only if the agent's deepest concerns lie beyond her own well-being. If they do, some sacrifices serve them and cohere with what she most cares about; if they do not, every sacrifice runs against all of them. Whether a sacrifice is rational in this sense is distinct from whether there is reason for it, and from whether it is admirable.

\begin{thebibliography}{99}\small
\bibitem[Frankfurt(1982)]{Frankfurt1982} Frankfurt, H.~G. (1982). The importance of what we care about. \emph{Synthese}, 53(2), 257--272.
\bibitem[Heathwood(2011)]{Heathwood2011} Heathwood, C. (2011). Preferentism and self-sacrifice. \emph{Pacific Philosophical Quarterly}, 92(1), 18--38.
\bibitem[Overvold(1980)]{Overvold1980} Overvold, M.~C. (1980). Self-interest and the concept of self-sacrifice. \emph{Canadian Journal of Philosophy}, 10(1), 105--118.
\bibitem[Parfit(1984)]{Parfit1984} Parfit, D. (1984). \emph{Reasons and Persons}. Oxford University Press.
\bibitem[Portmore(2007)]{Portmore2007} Portmore, D.~W. (2007). Welfare, achievement, and self-sacrifice. \emph{Journal of Ethics and Social Philosophy}, 2(2).
\bibitem[Sidgwick(1874)]{Sidgwick1874} Sidgwick, H. (1874). \emph{The Methods of Ethics}. London: Macmillan.
\end{thebibliography}

\end{document}
